Contact-rich manipulation requires actions to adapt as contact develops. During USB insertion, for example, the approach direction may remain useful over nearby control cycles, while local alignment must respond to newly observed contact. Force-aware Vision-Language-Action (VLA) policies incorporate interaction feedback into action generation~\cite{yu2025forcevlaenhancingvlamodels,li2026forcevla2unleashinghybridforceposition,zhang2025tavlaelucidatingdesignspace}, but typically predict task motion and contact-dependent adjustment together. Separating a reference action predicted by a slow pathway from a rapidly updated correction would allow contact feedback to affect execution without waiting for a new complete action chunk (Fig.~\ref{fig:teaser}).

Learning this correction from offline demonstrations is not straightforward. Demonstrations record the final action, without specifying which part belongs to the base behavior and which part is a contact-dependent adjustment. Multiple additive decompositions can reconstruct the same label. A correction pathway therefore needs a defined supervision target in addition to a fast architecture.

Existing asynchronous policies separate computation or sensor-update rates~\cite{li2026favlaforceadaptivefastslowvla,wang2026lateforceacceleratingvla,vanjani2026damvladecoupledasynchronousmultimodal}. Residual reinforcement learning and corrective imitation obtain correction signals from rewards or extra on-policy interaction~\cite{ankile2024imitationrefinementresidual,xu2025compliantresidualdaggerimproving}. We address the supervision problem using only the original offline demonstration dataset, while retaining asynchronous execution.

We present ForceDelta-VLA, which uses a lightweight policy to update force and delay corrections independently of reference-action generation. In offline training, paired force-conditioned and learned force-agnostic predictions of a frozen teacher define the force-correction target. The paired queries share task context, robot state, and sampling noise. A separate delay-correction target compensates for the mismatch between the previous reference action and the current force-agnostic prediction, including their different state references. Asynchronous schedule replay constructs these targets using the cached context available at deployment. During execution, the lightweight policy adjusts reference actions from the slow pathway using recent observations, responding to changes in robot state and contact between reference-action updates. The paired teacher queries are used only to construct training targets.

We evaluate ForceDelta-VLA on five single-arm and four bimanual real-robot tasks spanning contact establishment, precision insertion, friction-dominated motion, and constrained manipulation. The method achieves an $82.2\%$ mean success rate, compared with $54.4\%$ for ForceVLA and $70.6\%$ for direct rollout of the Temporal Teacher. Relative to ForceVLA, it reduces mean successful-trial peak contact force by approximately $26\%$ on both platforms. The correction forward pass takes $2.43$~ms.

The contributions of this work are:
\begin{enumerate}
    \item We introduce \textbf{force-agnostic correction distillation}: an explicit teacher-defined force-correction target from matched force-conditioned and learned force-agnostic predictions of a frozen teacher, using only the original offline dataset;
    \item We formulate a \textbf{state-consistent dual-correction decomposition} that separately supervises force and delay corrections in a common action coordinate system, with the delay target incorporating both reference-action mismatch and state rebasing; and
    \item We develop an \textbf{independently executable correction pathway} that updates pose corrections from recent force history and robot state while reusing reference actions and task context from the slow pathway.
\end{enumerate}
